\section{Integration back-ends, tolerance and number formats}\label{app:backends}

This appendix collects the detailed measurements behind the statements of
Sections~\ref{sec:method}--\ref{sec:parallel} on the integration back-ends
(the phase ODE, adaptive quadrature and a fixed-step trapezoid), on the
solver tolerance, and on narrow number formats. The integration back-ends
were measured on an earlier machine than the discrete march
(Section~\ref{sec:performance-unwrap}); their ratios carry over, their
absolute chart times do not.

\subsection{Convergence and efficiency of the integrators}\label{sec:performance-convergence}

\begin{figure}[tb]
  \centering
  \maybefig{fig_convergence}{\linewidth}
  \caption{Integration back-ends at an unstable point of the
  showcase system, where the exact count $\Zint=2$ is known (verified
  independently by a high-resolution semi-discretization eigenvalue count),
  so the plotted count error $|\Zraw-\Zint|$ is an \emph{absolute} error
  measure requiring no reference. (a)~Accuracy versus requested tolerance;
  the dotted line marks the practical rule
  $\mathrm{error}\approx100\,\mathrm{tol}$. (b)~Cost versus
  tolerance. (c)~Efficiency front. Error is on the vertical axis in both (a)
  and (c). A curve ends where a single parameter point exceeded a cost cap of
  one second, which is itself informative.}\label{fig:convergence}
\end{figure}

\begin{table}[tb]
  \centering
  \caption{Count error $|\Zraw-\Zint|$ (left block) and single-point CPU time
  (right block) of the back-ends at selected tolerances; system and point as
  in Fig.~\ref{fig:convergence}.}\label{tab:convergence}
  \resizebox{\linewidth}{!}{\maybeinput{tables/tab_convergence.tex}}
\end{table}

Figure~\ref{fig:convergence} and Table~\ref{tab:convergence} quantify the
choice of back-end, and three points deserve emphasis. First, from the
default tolerance downwards the adaptive back-ends track the requested
tolerance -- with the error settling roughly two orders of magnitude above
it, the rule of thumb~\eqref{eq:rule} of Section~\ref{sec:showcase-tolerance}
-- until they meet a common floor of $\approx8\cdot10^{-7}$ set by the tail
truncated at $\wmax=10^{6}$ and the estimated leading order, not by any
integrator; at loose tolerances the lower-order pairs do not yet track it
(Table~\ref{tab:convergence}). Second, the order of the pair matters far more
than its family.
Across the loose half of the ladder the classical fifth-order
pair~\cite{tsitouras2011runge} under-resolves the $\omega^{-1}$ ripple of
this velocity-feedback system, leaving count errors of $0.2$--$0.34$ --
within a hair of the rounding threshold -- and even at the package default
$10^{-5}$ it leaves $9\cdot10^{-2}$ where the ninth-order
default~\cite{verner2010numerically} leaves $6\cdot10^{-5}$: three orders
of magnitude, for a factor of $2.4$ in time. The linearly implicit
Rosenbrock
pair~\cite{rosenbrock1963some,shampine1997matlab,hairer1996solving} is
dominated from the default tolerance downwards, as
expected once one recognizes that the problem has no stiffness for it to
exploit (Section~\ref{sec:phaseode}); at tight tolerances it needs seconds
for a single parameter point where the ninth-order pair needs a quarter of
one, which is why the sweep in Fig.~\ref{fig:convergence} stops each method
once a point exceeds the cost cap. Third, at the default setting every
back-end still rounds to the correct count, so the choice of integrator
does not change the chart -- what differs, by orders of magnitude, is the
\emph{margin} to the rounding threshold, and the default is chosen for that
margin rather than for raw speed.


\subsection{Tolerance study on the showcase system}\label{app:tolerance}


\begin{figure}[tb]
  \centering
  \maybefig{fig_tolerance_charts}{\linewidth}
  \caption{The showcase chart computed with solver tolerances
  $10^{-2}$, $10^{-4}$ and $10^{-6}$ (left to right); the
  wall-clock time of each full sweep is given in the panel titles. The charts
  are visually near-identical even at $10^{-2}$, where the count is already
  uncertain: the tolerance buys certainty rather than appearance, cf.\
  Fig.~\ref{fig:errorfield} and Table~\ref{tab:tolerance}. The study uses
  the standard chart window $\wmax=10^{4}$ and the reduced
  form~\eqref{eq:showcasered} of the characteristic
  function.}\label{fig:tolerance}
\end{figure}

\begin{figure}[tb]
  \centering
  \maybefig{fig_error_field}{\linewidth}
  \caption{Self-validating error: fields and histograms of the integer
  residual $\varepsilon=|\Zraw-\operatorname{round}(\Zraw)|$,
  cf.~\eqref{eq:eps}, for the three tolerances of
  Fig.~\ref{fig:tolerance}. Dashed red line: the requested tolerance; dotted
  black line: the practical rule $\varepsilon\approx100\,\mathrm{tol}$. The
  residual scales with the tolerance uniformly over the parameter plane,
  amplified by the accepted-step count of the frequency march, until the
  tail truncated at $\wmax=10^{4}$ (of order $2\cdot10^{-5}$) becomes the
  dominant error.}\label{fig:errorfield}
\end{figure}

\begin{table}[tb]
  \centering
  \caption{Tolerance study on the showcase system
  (grid as in Fig.~\ref{fig:tolerance}): sweep time, integer residual
  statistics, the number of grid points where the rounded count would be
  ambiguous ($\varepsilon\ge1/4$), the number whose $\Zint$ actually deviates
  from the $10^{-6}$-tolerance reference, and the subset of those that flip
  the stable/unstable classification. The ``uncertain'' and ``wrong'' columns
  measure genuinely different things: a skipped peak shifts $\Zraw$ by
  exactly one and leaves it just as close to an integer, so a point can be
  wrong while looking certain -- which is precisely why the
  cross-check~\eqref{eq:crosscheck} is needed.}\label{tab:tolerance}
  \maybeinput{tables/tab_tolerance.tex}
\end{table}

Figure~\ref{fig:tolerance} shows the identical chart computed with
tolerances $10^{-2}$, $10^{-4}$ and $10^{-6}$, together with the
measured sweep times, and Fig.~\ref{fig:errorfield} shows the corresponding
integer-residual fields~\eqref{eq:eps}; the ladder deliberately starts
beyond the safe range. Three observations matter for practice. First, the
tolerance buys \emph{certainty}, not appearance. The panels are visually
near-identical, yet the loosest one is genuinely broken: measured against
the $10^{-6}$ reference, $10^{-2}$ miscounts $349$ of the $8000$ points and
flips the stable/unstable classification on $89$
(Table~\ref{tab:tolerance}). That failure is anything but silent -- the
residual field of Fig.~\ref{fig:errorfield} lights up across the whole
chart, with $210$ points beyond the rounding threshold -- and one decade
into the safe range the problem is essentially gone: at $10^{-4}$ a single
count of the $8000$ deviates, by one and without flipping the
classification, and at $10^{-6}$ none. Second, while the integrator
dominates the error budget, the largest integer residual of the chart tracks
the requested tolerance with a roughly constant factor,
$\max\varepsilon\approx100\,\mathrm{tol}$, Eq.~\eqref{eq:rule}, to within a factor of four (Table~\ref{tab:tolerance}; the mean residual is
one to two orders of magnitude smaller); the factor is inherited from the
accepted-step count of the march. At the standard window $\wmax=10^{4}$ the
truncated tail of this velocity-feedback system, $\approx2\cdot10^{-5}$,
takes over once the tolerance drops below $\sim10^{-6}$ -- the mean
residual of the tightest panel settles exactly there -- and, because every
error source surfaces in the same number, the diagnostic itself tells the
user that the next decade of certainty must come from $\wmax$, not from the
tolerance. The accuracy of $\Zint$ can therefore be \emph{prescribed}
rather than discovered: a user who wants the count certain to $10^{-3}$
asks for $10^{-5}$, which is exactly the package default. Third, the
CPU-time premium of tightening by four decades is a factor of three
(Table~\ref{tab:tolerance}), because the extra effort is spent only near
the integrand peaks.

For chart work a miscount is a milder error than a displaced boundary, and
the two degrade differently: the count is a rounded quantity and fails
discontinuously, whereas the tracked $\sest$ inherits the tolerance
smoothly. The traced boundary is built from $\sest$ alone through the
objective~\eqref{eq:objective}, and it stays a small fraction of one grid
pixel away from the reference trace at every tolerance studied
(\code{data/boundary\_shift.csv}). The practical recipe follows: a
moderately loose tolerance ($10^{-3}$--$10^{-4}$) is an appropriate tool
for an initial search -- the residual field will announce if it was too
loose -- after which the interesting region is recomputed at the default
setting to certify the counts, with the cross-check~\eqref{eq:crosscheck}
guarding the one error mode near the boundary that the residual cannot
reveal.

A flagged pixel need not be recomputed with a different back-end: the march
can repair itself (\code{peak\_repair=true}). The skipped peak is easy to
\emph{locate} even when it is far too narrow to \emph{integrate}: its center
is a minimum of $|\Dfun(\sigma+\ii\omega)|^2$, which dips toward that
minimum over an $O(1)$ frequency range no matter how narrow the peak is. A
root-finding callback stops the integrator at every such minimum,
recomputes the phase increment of the crossing step by quadrature whose
subdivision is forced into the located peak, and restarts the march at the
minimum. Measured across the Hopf point of this chart
(\code{data/peak\_repair.csv}): the plain march loses one count from
$\Delta P = \RepairPlainFailDP$ inward, the Gauss--Kronrod back-end -- which
must \emph{discover} the peak -- loses it silently at
$\Delta P = \RepairDeepDP$, whereas the repaired march, which \emph{knows}
where the peak is, still counts correctly there. The repair matters only in
the immediate vicinity of the boundary, and arming it for a whole sweep
costs $\approx\RepairCostFactor\times$ the plain march, which is why it is
off by default; the economical pattern is to sweep plainly, flag
with~\eqref{eq:crosscheck}, and re-evaluate only the flagged pixels with
the repair on.


\subsection{Chart-level cost, memory and cold start}\label{sec:performance-cost}

\begin{table}[tb]
  \centering
  \caption{Warm chart-level costs (median of repeated runs,
  \NumThreads{}~threads): full brute-force sweeps for three systems of
  increasing evaluation cost, with allocated memory. The last column counts
  the points at which the back-end's $\Zint$ deviates from a
  tight-tolerance reference of the same chart: the adaptive back-ends are
  required to match everywhere before their time is recorded, while the
  fixed-step back-end -- whose resolution is set by hand -- reports its
  misses.}\label{tab:gridtimings}
  \resizebox{\linewidth}{!}{\maybeinput{tables/tab_grid_timings.tex}}
\end{table}

\begin{table}[tb]
  \centering
  \caption{The two operating points of the method on the same
  $100\times100$ chart: the \emph{accurate} preset (the default tolerance)
  and a \emph{fast} preset for first exploration. The last column counts
  the points whose $\Zint$ differs from a tight reference, i.e.\ the price
  of the fast preset, which buys roughly a factor of two in
  time.}\label{tab:presets}
  \maybeinput{tables/tab_presets.tex}
\end{table}

Two operating points matter in practice, and Table~\ref{tab:presets}
separates them. With the \emph{accurate} preset -- the standard chart
window $\wmax=10^{4}$ and the default $\mathrm{tol}=10^{-5}$ -- a complete
$100\times100$ chart of a scalar characteristic function costs a fraction
of a second, and every one of the ten thousand counts matches a tight
reference. A \emph{fast} preset ($\mathrm{tol}=10^{-3}$) buys roughly a
factor of two at the price of a few dozen boundary-adjacent miscounts --
acceptable for a first look, and flagged by the diagnostics of
Section~\ref{sec:showcase-tolerance}. Beyond the presets, two user choices
dominate the cost on the showcase DDAE, spanning three orders of magnitude
together: $\wmax$ (the velocity feedback's $\omega^{-1}$ ripple makes
$\wmax=10^{6}$ cost some fifty times $\wmax=10^{4}$, while buying nothing
-- the truncated tail is five orders below the rounding threshold) and the
form of $\Dfun$ (the automatic extraction costs a factor of $\ExtractRatio$
per evaluation over the hand-derived expression,
Table~\ref{tab:extraction}). Choosing badly on both turns the $1.98$\,s
chart of Table~\ref{tab:presets} into one of more than half an hour; this is what the
asymmetric guidance of Section~\ref{sec:phaseode} means in practice.

Table~\ref{tab:gridtimings} reports the cost of a \emph{complete}
$100\times100$ chart -- ten thousand parameter points -- on a
\CPUmodel{} (\NumThreads{} threads, \RAMgb\,GB RAM, Julia~\JuliaVersion,
package revision \PkgCommit), with the per-point cost listed alongside so
that other resolutions can be scaled directly. For scalar characteristic
functions a full chart takes tens to hundreds of milliseconds, and the
entire hybrid workflow (colored sweep plus MDBM boundary) finishes well
inside one second -- genuinely interactive. The last column is worth a
glance: the fixed-step back-end is the fastest entry for the delayed systems
and gets thousands of the ten thousand counts wrong, for the reason given in
Section~\ref{sec:alternatives}. The showcase DDAE is listed with its
hand-derived characteristic function, as are all figures in this paper; what
the automatic extraction of Section~\ref{sec:extraction} would cost instead
is measured separately in Table~\ref{tab:extraction}. Memory stays modest throughout
because the phase ODE carries a single scalar state and the sweep stores
only the five result fields; no large matrices are formed at any point.
Cold start is a one-time cost dominated by package loading: in a fresh Julia
process, loading takes $5$--$8$\,s and the first thousand-point sweep
$0.07$--$0.2$\,s thanks to the precompilation workload shipped with the
package, after which every further chart runs at the warm speeds of
Table~\ref{tab:gridtimings}.

\subsection{A broader solver comparison}\label{sec:performance-zoo}

\begin{figure}[tb]
  \centering
  \maybefig{fig_solver_zoo}{0.7\linewidth}
  \caption{Error--cost Pareto fronts of eight integration back-ends on a
  $40\times40$ sweep of the fourth-order benchmark oscillator
  (\ref{app:gallery}): fixed-step trapezoid (step-count sweep),
  adaptive Gauss--Kronrod, and six ODE solvers (tolerance
  sweeps).}\label{fig:zoo}
\end{figure}

\begin{table}[tb]
  \centering
  \caption{Evaluations of $\Dfun$ needed for one parameter point at
  tol~$=10^{-5}$, with the count of each entry verified against a
  tight-tolerance reference before being recorded. What governs the cost is
  how fast the delay-induced ripple of $\theta'$ decays, i.e.\ how far the
  delayed terms sit below the leading power: $\omega^{-1}$ for the
  velocity-feedback showcase, $\omega^{-3}$ for the fourth-order
  position-feedback benchmark.}\label{tab:evals}
  \maybeinput{tables/tab_evals.tex}
\end{table}

Table~\ref{tab:evals} isolates what actually drives the cost. The decisive
quantity is not the back-end but the decay rate of the phase ripple, set by
how far the delayed terms sit below the leading power of $\Dfun$. Where the
ripple decays fast (delayed \emph{position} feedback,
$\theta'\sim\omega^{-3}$), quadrature recognizes the negligible tail panel
by panel and finishes about $2.5$ times cheaper than the march; where it
decays slowly (delayed \emph{velocity} feedback, $\theta'\sim\omega^{-1}$),
the ranking reverses and the march is some three times cheaper. Neither
back-end dominates, and the lesson for a user is to spend attention on
$\wmax$ and on the structure of the feedback, not on the choice of
integrator. One implementation detail underpins the quadrature numbers: the
Gauss--Kronrod back-end is seeded with logarithmically spaced interior
breakpoints, one per two frequency decades. Without them the first
$G_7$--$K_{15}$ panel spans the entire window, its lowest node lands above
the whole resonance region for $\wmax=10^{6}$, and the routine returns
$\Zraw=n/2$ having never seen a single root -- with an integer residual of
\emph{exactly} zero, since nothing about the quadrature failed. The
breakpoints cost a few dozen extra evaluations and remove the failure mode
entirely.

Figure~\ref{fig:zoo} extends the comparison to a wider selection of ODE
solvers on the fourth-order benchmark system. The fixed-step trapezoid is
\emph{dominated everywhere}: even at $3\cdot10^{4}$ steps -- more than an
order of magnitude slower than the adaptive default -- its worst error over
the chart remains several roots, because a uniform grid cannot resolve a
peak that narrows without bound as a root approaches the line. All adaptive
methods follow their tolerance; higher-order pairs dominate at tight
tolerances, and the linearly implicit pair is dominated. The spread between
the \emph{adaptive}
methods is far smaller than the gap that separates all of them from a fixed
grid, which is the practical message of this section.
\subsection{Narrow number formats}\label{app:precision}

\begin{table}[tb]
  \centering
  \caption{The discrete march with $\Dfun$ evaluated in narrow number formats
  ($400\times400$ chart of the fourth-order benchmark, $\wmax=15$, best of four
  frequency scalings $\lam=\Omega\mu$ per format; the march itself runs in
  Float32). \emph{Flagged} points carry a sub-resolution or residual flag and
  are re-checked in a wider format; \emph{failed} points did not reach $\wmax$.
  $^\ast$Emulated: every operation is rounded to the format.}\label{tab:precision}
  \maybeinput{tables/tab_precision.tex}
\end{table}

Since the march needs only the
\emph{phase} of $\Dfun$ to a fraction of a radian, the characteristic function
can even be evaluated in half precision (Table~\ref{tab:precision}). Float16
is accurate wherever no root is close to the line: after rescaling the
frequency into its range it yields, on the full-HD chart of
Table~\ref{tab:gpu}, $\GpuFsixteenWrong\,\%$ wrong counts that are \emph{all}
flagged, and $\GpuFsixteenFlagged\,\%$ flagged points in total; re-checking the
flagged points in Float32 -- in place on the device: gather, march, scatter --
restores the exact chart. It does not, however, make the chart faster: at 8K
the Float16 evaluation takes $1.01$--$1.12$ times the Float32 time on all four
cards (the factor of three of the T4 at full HD is an artifact of its clock
control). Scalar half-precision instructions issue at the single-precision
rate -- a gain would require packing pairs of values into vector registers --
the elementary functions are evaluated by promotion to Float32, and the
conversions have a reduced throughput of their own. (An earlier version of the
kernel, which passed the dual numbers through local memory, did run three
times faster in Float16, \code{gpu/results/interactive\_G4.csv}; the gain came
from the smaller memory traffic and disappeared once that traffic was
removed.) Float32 is therefore the production format on every card measured,
and the value of the narrow formats lies in what they reveal about the method:
how little precision the phase needs. BFloat16, with only seven mantissa bits,
flags $14\,\%$ of the chart, and the 8-bit and 4-bit formats are useless for
this purpose: with one to three mantissa bits even the \emph{parameters} of a
chart are quantized to a handful of distinct values per octave, so that
$95$--$100\,\%$ of the points are flagged and up to one in ten counts is
wrong; Float4 gets every count wrong. Load balancing, finally, does not pay:
the dynamic work queue of Fig.~\ref{fig:parallel}, in which threads that finish
early fetch new points, is $1.3$--$1.9$ times \emph{slower} than one thread per
point on every card, because the cost per point varies only moderately and the
hardware scheduler absorbs the difference at a lower overhead than the queue.
